% ======================================================================
% main.tex
% Triple Humanoid 24/7 Adverse Event Oncology Trial Response Team:
% 4-Site Rotation
%
% Author: Kevin Kawchak, CEO ChemicalQDevice
% Orcid: 0009-0007-5457-8667
% DOI placeholder: 10.5281/zenodo.xxxxxxxx
% Date: May 20, 2026
%
% This file is a draft scaffold. Every section file contains bracketed
% instructions that point a future Claude Code Opus 4.7 1M Max session
% at the exact directories and file names it should read in order to
% replace each placeholder block with finished paper prose. The future
% session is expected to expand this scaffold to a 70 plus page paper.
% No section in this scaffold should be treated as final text. Read the
% bracketed instructions at the top and inside each section file, then
% read the files those instructions point to, then write the section.
% Do not modify any file under paper/codegen, paper/execution,
% paper/instructions, or paper/inputs while expanding this scaffold.
% Those four trees are read only RESEARCH inputs.
% ======================================================================
\documentclass[11pt]{article}
\usepackage{humanoid_paper_template}

\title{Triple Humanoid 24/7 Adverse Event Oncology Trial Response Team: 4-Site Rotation}
\author{%
  Kevin Kawchak\hspace{2pt}\orcidicon{0009-0007-5457-8667}\\
  CEO ChemicalQDevice\\
  \doiLink{10.5281/zenodo.xxxxxxxx}
}
\date{May 20, 2026}

\begin{document}
\maketitle

\begin{center}
\small\itshape
Disclaimer: This work is independent and not endorsed or sponsored by trial sponsors, FDA, CRO, site, IRB, regulator, or medical society; and was generated using Artificial Intelligence.
\end{center}

\vspace{0.4em}

\begin{abstract}
% ======================================================================
% sections/abstract.tex
%
% Abstract scaffold for the paper:
%   Triple Humanoid 24/7 Adverse Event Oncology Trial Response Team:
%   4-Site Rotation
%
% This file is a draft scaffold. A future Claude Code Opus 4.7 1M Max
% session reads the bracketed instruction blocks below, then reads the
% files those instructions name, then expands the scaffold into the
% finished abstract text. The future session should overwrite both the
% bracketed instructions and the placeholder prose with the finished
% paragraph, while keeping the abstract length to roughly 300 words.
% ======================================================================

\begin{sectionaccent}

[ABSTRACT INSTRUCTIONS FOR FUTURE CLAUDE CODE OPUS 4.7 1M MAX SESSION.
The abstract has two halves and reads as a single continuous paragraph
of about 300 words. The first half is the production overview: it must
walk the reader through the full chain of code instructions, code
generations, code execution, draft paper generation, and full paper
production. The second half is the results catch: it must surface
specific numerical results, discussion claims, and metrics that appear
elsewhere in the paper so that a reader who only reads the abstract
still walks away with concrete numbers. Do not over generalize. Use
the file list below.

Files to read before writing the production overview half:
\begin{itemize}
  \item \path{demo-projects/07-humanoid/paper/instructions/README.md}
        for the project thesis, the four site identifiers (SF-01, SD-01,
        BO-01, AT-01), the 12 H2 robot total count, the 1 Hz LLM
        broadcast cadence, the 10 Hz humanoid motion cadence, the 168
        hour monitoring window, and the camaraderie reduction factor of
        3 in single robot error potential.
  \item \path{demo-projects/07-humanoid/paper/instructions/00-project-overview/}
        through \path{09-runtime-instructions/} for the 10 instruction
        directories that drive code generation.
  \item \path{demo-projects/07-humanoid/paper/codegen/README.md} and
        \path{demo-projects/07-humanoid/paper/codegen/architecture.md}
        for the code generation phase output.
  \item \path{demo-projects/07-humanoid/paper/execution/README.md}
        and \path{demo-projects/07-humanoid/paper/execution/reports/execution_report_v0_5_0.md}
        for the execution phase output.
  \item \path{demo-projects/07-humanoid/paper/inputs/README.md},
        \path{demo-projects/07-humanoid/paper/inputs/Deep-Research-A/README.md},
        and the equivalent README in
        \path{demo-projects/07-humanoid/paper/inputs/Deep-Research-B/}
        for the prior paper context (the prior paper covered Medical
        Full Automation A and Medical Rapid Prototyping B).
\end{itemize}

Files to read before writing the results catch half:
\begin{itemize}
  \item \path{demo-projects/07-humanoid/paper/codegen/data/iterations/index.jsonl}
        for the 32 deterministic iteration sweep results across the 4
        sites and many columns. Quote the iteration count and the
        per-site count in the abstract.
  \item \path{demo-projects/07-humanoid/paper/codegen/reports/dashboard.html}
        and \path{demo-projects/07-humanoid/paper/codegen/reports/report.md}
        for the comparison summary headline numbers. Surface the lead
        metric (the camaraderie factor of 3, plus the e-stop latency of
        5 ms swarm wide).
  \item \path{demo-projects/07-humanoid/paper/execution/data/site_runtime_runs/AT-01_llm_decisions.jsonl}
        plus \texttt{BO-01}, \texttt{SD-01}, and \texttt{SF-01} variants
        in the same directory for the per-site LLM decision streams at
        1 Hz. Note that the 4 site files are identical in byte length
        (62,330 bytes each); call this out as a check that the runtime
        produced parallel-scope decision logs across sites and not as a
        data duplication artifact in production.
  \item \path{demo-projects/07-humanoid/paper/execution/data/iterations_aggregate.duckdb}
        for the DuckDB roll up of all 32 iteration sweeps.
\end{itemize}

Required content cues for the abstract paragraph:
\begin{itemize}
  \item Open with the thesis sentence verbatim from the paper README.
  \item Name the 4 sites (SF-01, SD-01, BO-01, AT-01), the 3 H2 per
        site, the 12 total H2, and the 1 on-prem Claude Opus 4.7 1M
        instance per site.
  \item Name the 1 Hz broadcast cadence and the 10 Hz robot motion
        cadence.
  \item Name the 168 hour monitoring window and the approximate 84 AE
        plus 24 SAE event volume.
  \item Name CTCAE version 5 and the FDA RTCT 1 hour SLA from HR 9505
        as the regulatory anchor.
  \item Name the 32 iteration sweeps and the camaraderie reduction
        factor of 3 in single robot error potential.
  \item Cite the author's December 2025 paper
        (\path{kawchak_2025_18029100}) as the prior evidence that
        LLMs are effective on FDA adverse events data.
  \item Close on the future industry value: the on-prem repository LLM
        plus sensor pipeline plus Cartesian x, y, z command structure
        is presented as a checklist already executed so the industry
        does not have to repeat it.
\end{itemize}

Do not exceed roughly 300 words. Do not introduce numbers that are not
supported by the files named above. Do not use em dashes, double
dashes, or triple dashes. Replace any stray \texttt{SS} with the
section sign \S\ during the cleanup pass.]

The abstract will summarize the on-premises repository based LLM
pipeline, the 4-site rotation across SF-01, SD-01, BO-01, and AT-01,
the 3 Unitree H2 EDU robots per site, the 1 Hz broadcast cadence, the
Cartesian x, y, z command surface, and the per-site camaraderie
behavior that reduces single robot error potential. The closing
sentences will surface specific results from the 32 iteration sweeps
and the per-site LLM decision streams.

\sectiontable{Abstract planning matrix}{%
Thesis & A1 & 1 & instructions README & State thesis verbatim & 1 sentence & Reads like the README & Ready \\
Overview & A2 & 2 & code, exec README & Production chain & 4 to 5 sent. & Names phases & Ready \\
Numbers & A3 & 3 & index, dashboard & Headline metrics & 3 to 4 sent. & Sourced numbers & Ready \\
}

\end{sectionaccent}

\end{abstract}

\noindent\textbf{\keywordsname:}\ physical AI; humanoid robotics; Unitree H2; on-premises large language model; Cartesian control; clinical trial adverse events; CTCAE; FDA RTCT pilot; oncology; swarm camaraderie; 4-site rotation.

\section{Introduction}
\label{sec:introduction}
% ======================================================================
% sections/introduction.tex
%
% Introduction scaffold for the paper:
%   Triple Humanoid 24/7 Adverse Event Oncology Trial Response Team:
%   4-Site Rotation
%
% This file is a draft scaffold. A future Claude Code Opus 4.7 1M Max
% session reads the bracketed instruction blocks below, then reads the
% files those instructions name, then expands the scaffold into the
% finished introduction. Target length is roughly 6 to 8 pages of the
% 70 plus page final paper.
% ======================================================================

\begin{sectionaccent}

[INTRODUCTION INSTRUCTIONS FOR FUTURE CLAUDE CODE OPUS 4.7 1M MAX SESSION.
Write the introduction as 5 connected subsections inside this single
\texttt{\textbackslash section\{Introduction\}} block. The 5 subsections are
1) Background And Setting, 2) Thesis Statement And Scope, 3) The
On-Premises Repository LLM Pattern, 4) The 4-Site Continental Rotation,
and 5) Roadmap Of The Paper. Use the bracketed file lists below to
draw the content. Add a verbatim ASCII diagram lifted from the file
named below into the 4-Site Continental Rotation subsection.

Files to read before writing Background And Setting:
\begin{itemize}
  \item \path{demo-projects/07-humanoid/paper/instructions/README.md}
        (whole file).
  \item \path{demo-projects/07-humanoid/paper/instructions/00-project-overview/}
        (whole directory, every file).
  \item \path{demo-projects/07-humanoid/paper/inputs/Deep-Research-A/chunk1_paper_text.md}
        and \path{chunk2_tables_references.md} for the Medical Full
        Automation A context.
  \item \path{demo-projects/07-humanoid/paper/inputs/Deep-Research-B/chunk_01_executive_summary_and_fully_automated_sponsor.md},
        \path{chunk_02_federated_learning_and_patient_prediction.md},
        \path{chunk_03_clinagent.md}, and
        \path{chunk_04_comparative_overview_and_implications.md}
        for the Medical Rapid Prototyping B context.
  \item Author's December 2025 paper
        \path{kawchak_2025_18029100} for the prior evidence that
        LLMs handle FDA adverse event data effectively. Cite the bib
        key \path{kawchak_2025_18029100} and quote the DOI
        \path{10.5281/zenodo.18029100}.
\end{itemize}

Files to read before writing Thesis Statement And Scope:
\begin{itemize}
  \item \path{demo-projects/07-humanoid/paper/instructions/README.md}
        Thesis section (verbatim quote of the thesis sentence is
        required in this subsection).
  \item \path{demo-projects/07-humanoid/paper/codegen/architecture.md}
        whole file.
\end{itemize}

Files to read before writing The On-Premises Repository LLM Pattern:
\begin{itemize}
  \item \path{demo-projects/07-humanoid/paper/codegen/src/llm_planner.py}
        for the 1 Hz planner loop.
  \item \path{demo-projects/07-humanoid/paper/codegen/src/broadcaster.py}
        for the fan out structure that delivers the same command
        payload to all 3 robots at a site at the same tick.
  \item \path{demo-projects/07-humanoid/paper/codegen/src/prompt_frozen.md}
        for the frozen prompt content.
  \item \path{demo-projects/07-humanoid/paper/codegen/src/comms_intellectual.py}
        and
        \path{demo-projects/07-humanoid/paper/codegen/src/comms_physical.py}
        for the dual channel (60 GHz UWB plus IR beacon plus shared
        compute fabric) that turns the per-site Claude Opus 4.7 1M
        instance into the on-premises repository LLM.
\end{itemize}

Files to read before writing The 4-Site Continental Rotation:
\begin{itemize}
  \item \path{demo-projects/07-humanoid/paper/codegen/diagrams/network_layout.txt}
        Include this ASCII diagram verbatim inside an \texttt{asciibox}
        environment as Figure 1. Caption it
        ``PAT-NET-001 continental layout with one Claude Opus 4.7 1M
        on-prem broadcaster per site and 3 Unitree H2 EDU robots per
        site.''
  \item \path{demo-projects/07-humanoid/paper/codegen/diagrams/swarm_dance.txt}
        Brief callout reference. Save the full inclusion for the
        Methods section.
  \item \path{demo-projects/07-humanoid/paper/codegen/src/cross_site_bus.py}
        for the de-identified cross site summary bus.
  \item \path{demo-projects/07-humanoid/paper/execution/diagrams/02_4site_execution_sequence.txt}
        for the 4 site execution sequence. Include this ASCII diagram
        verbatim inside an \texttt{asciibox} environment as Figure 2.
\end{itemize}

Files to read before writing Roadmap Of The Paper:
\begin{itemize}
  \item \path{demo-projects/07-humanoid/paper/draft-paper/main.tex}
        for the section ordering (Methods, Results, Discussion,
        Limitations and Future Work, Conclusions).
\end{itemize}

Required content cues for the introduction:
\begin{itemize}
  \item Open by framing the United States obligation to lead in patient
        safety, efficacy, and speed for oncology adverse event
        response. Stay respectful but advocate for accelerating the
        new field of Physical AI oncology trial surgical procedures.
  \item Frame the FDA RTCT 1 hour SLA as a real, accountable timer.
        Pair it with the camaraderie factor of 3 reduction in single
        robot error.
  \item Cite \path{kawchak_2025_18029100} once when introducing the
        FDA adverse event reporting evidence base.
  \item Present each milestone in the production chain as already
        checked off so the industry can build on top of it rather
        than rebuilding it from scratch.
  \item Use the ASCII diagrams listed above verbatim. Do not redraw.
\end{itemize}

Required real-life feasibility insight to include in this section:
``As of May 2026 the Unitree H2 EDU is shipping at the spec sheet level
named in the instructions README (39 DOF, 1.8 m, 70 kg, IP65 outdoor
rated). Pairing it with one on-prem Claude Opus 4.7 1M instance per
site is feasible on conventional high end x86 or Apple Silicon
servers. The 1 Hz broadcast cadence is well within current network and
inference latency budgets. What remains is the IRB, sponsor, and FDA
RTCT pilot enrollment, all of which the paper treats as part of the
forward roadmap rather than as solved.''

Do not exceed 8 pages. Use single dashes only. Use \S\ for the section
sign. Avoid stranding 1 or 2 word lines on their own.]

The introduction will set the FDA RTCT pilot context, name the 4-site
continental network, restate the thesis, and walk the reader through
the on-premises repository LLM pattern, the per-site swarm camaraderie,
and the roadmap that the rest of the paper follows.

\sectiontable{Introduction planning matrix}{%
Context & I1 & 1 & instr, inputs & FDA RTCT & 1 to 2 paragraphs & Cites kawchak & Ready \\
Thesis & I2 & 2 & instructions README & Quote thesis sentence & 1 paragraph & Verbatim quote & Ready \\
Pattern & I3 & 3 & llm planner & On-prem LLM & 2 paragraphs & 1 Hz cadence & Ready \\
Network & I4 & 4 & net layout & ASCII Figure 1 & 1 figure & 4 sites visible & Ready \\
Roadmap & I5 & 5 & main.tex & Section walk & 1 paragraph & Reader path & Ready \\
}

\end{sectionaccent}


\newpage
\phantomsection
\addcontentsline{toc}{section}{Table of Contents}
\tableofcontents
\newpage

\section{Methods}
\label{sec:methods}
% ======================================================================
% sections/methods.tex
%
% Methods scaffold for the paper:
%   Triple Humanoid 24/7 Adverse Event Oncology Trial Response Team:
%   4-Site Rotation
%
% This file is a draft scaffold. A future Claude Code Opus 4.7 1M Max
% session reads the bracketed instruction blocks below, then reads the
% files those instructions name, then expands the scaffold into the
% finished methods section. Target length is roughly 18 to 22 pages of
% the 70 plus page final paper.
% ======================================================================

\begin{sectionaccent}

[METHODS INSTRUCTIONS FOR FUTURE CLAUDE CODE OPUS 4.7 1M MAX SESSION.
Write the methods section as 9 subsections inside this single
\texttt{\textbackslash section\{Methods\}} block. Read every file named
below before writing the subsection that uses it. Do not invent
implementation details. Use the file names verbatim in the prose so a
reader can navigate to the source code.

Subsection 1, Code Instruction Layer. Read:
\begin{itemize}
  \item \texttt{demo-projects/07-humanoid/paper/instructions/README.md}
        (whole file).
  \item Every \texttt{README.md} in the 10 numbered subdirectories of
        \texttt{demo-projects/07-humanoid/paper/instructions/}:
        \texttt{00-project-overview/}, \texttt{01-commit-roadmap/},
        \texttt{02-config-instructions/},
        \texttt{03-schemas-instructions/},
        \texttt{04-robotic-instructions/},
        \texttt{05-cartesian-instructions/},
        \texttt{06-iteration-instructions/},
        \texttt{07-llm-planner-instructions/},
        \texttt{08-comparison-instructions/}, and
        \texttt{09-runtime-instructions/}. The
        \texttt{10-repository-update-instructions/} directory is
        named in the instructions README repository structure block
        but is not yet present on disk; mention this gap explicitly
        as a documented future commit.
\end{itemize}

Subsection 2, Code Generation Layer And The 27 \texttt{src/} Files.
Read every file in
\texttt{demo-projects/07-humanoid/paper/codegen/src/} and discuss the
function of each. Group the discussion by function rather than by
filename:
\begin{itemize}
  \item Sensor to Cartesian pipeline: \texttt{sensor\_to\_xyz.py},
        \texttt{kinematics.yaml}, \texttt{cartesian\_planner.py},
        \texttt{xyz\_safety.py}. Discuss how raw sensor frames become
        x, y, z coordinate goals subject to per arm cumulative force
        budget caps of 15 N during chest compressions and 5 N
        otherwise, and to an inter robot minimum distance of 1.2 m at
        rest and 0.4 m during shared task handoff, and a cumulative
        cross robot force cap of 22 N during 3 robot patient
        transfer.
  \item Per site dispatch: \texttt{h2\_dispatcher.py}. This file
        exhibits per site H2 robot dispatch proficiency and is the
        single point where role assignment becomes a robot command.
        Call out the file by name in the prose.
  \item Swarm camaraderie: \texttt{swarm\_coordinator.py} at 191 lines.
        This file implements the camaraderie logic that animates the
        3 robot Unitree H2 EDU camarade swarm at one site. It is used
        by the LLM planner to author broadcasts and by the in-robot
        loop to update local behavior. Discuss that 191 lines is a
        compact length for non trivial 3 way coordination logic; the
        compactness is a direct consequence of using the per-site
        broadcaster pattern and the shared state tracker rather than
        a peer to peer consensus protocol.
  \item LLM planner and broadcast: \texttt{llm\_planner.py},
        \texttt{broadcaster.py}, \texttt{prompt\_frozen.md},
        \texttt{comms\_intellectual.py}, \texttt{comms\_physical.py},
        \texttt{cross\_site\_bus.py}, \texttt{peer\_state\_tracker.py}.
        Note the 1 Hz broadcast cadence and the redundant failover
        within the site.
  \item CTCAE grading and FDA submission:
        \texttt{ctcae\_grader.py},
        \texttt{fda\_rtct\_submitter.py} at 67 lines. The submitter
        submits CTCAE grade 2 or higher AE records to the FDA RTCT
        pilot intake API within 1 hour per the HR 9505 SLA, emits
        records matching the \texttt{fda\_rtct\_submission} schema,
        and uses ed25519 signing in production with a placeholder in
        simulation. Discuss that 67 lines is the right scale for a
        signed JSON submitter and that the small surface area is a
        regulatory feature, not a limitation.
  \item Physician escalation and error checking:
        \texttt{physician\_escalation.py},
        \texttt{check\_errors.py}. The error checker provides LLM
        credibility to FDA and other regulators by trapping schema
        and physics rule violations before broadcast.
  \item Site runtime and iteration sweep:
        \texttt{site\_runtime.py}, \texttt{iterate.py},
        \texttt{week\_runner.py}, \texttt{ingest.py},
        \texttt{compute.py}.
  \item Comparison and figures: \texttt{compare\_agent.py},
        \texttt{figures\_gen.py}.
  \item Native code paths: \texttt{robot\_loop.cpp} for the in-robot
        10 Hz loop, \texttt{runner.rs} and \texttt{runner\_main.rs}
        for the Rust runner.
\end{itemize}
List the function of every file. The 27 files in
\texttt{codegen/src/} taken together represent a processing feat
because they cover dispatch, planning, kinematics, comms, grading,
submission, error checking, iteration, runtime, comparison, figure
generation, native robot loop, and Rust runner inside one tree.

Subsection 3, Schemas And Configs. Read every file in
\texttt{demo-projects/07-humanoid/paper/codegen/schemas/} and
\texttt{demo-projects/07-humanoid/paper/codegen/config/}.

Subsection 4, Cartesian Coordinate Surface. Read
\texttt{demo-projects/07-humanoid/paper/codegen/src/kinematics.yaml},
\texttt{cartesian\_planner.py}, \texttt{sensor\_to\_xyz.py}, and
\texttt{xyz\_safety.py} again with a focus on the x, y, z command
representation. Discuss that the Cartesian surface is the contract
between the LLM and the robot; the LLM does not author joint angles.

Subsection 5, Camaraderie Behavior. Read the 191 line
\texttt{swarm\_coordinator.py} and the ASCII diagram
\texttt{demo-projects/07-humanoid/paper/codegen/diagrams/swarm\_dance.txt}.
Include the diagram verbatim inside an \texttt{asciibox} environment
as Figure 3.

Subsection 6, Iteration Sweep. Read
\texttt{demo-projects/07-humanoid/paper/codegen/data/iterations/index.jsonl}.
Quote the 32 iteration count across 4 trial sites and many columns as
direct evidence that Claude Code can create project code effectively.
Also read \texttt{demo-projects/07-humanoid/paper/execution/diagrams/03\_iteration\_sweep\_funnel.txt}
and include it verbatim inside an \texttt{asciibox} environment as
Figure 4.

Subsection 7, Tests And Error Checks. Read
\texttt{demo-projects/07-humanoid/paper/codegen/src/check\_errors.py}
and \texttt{demo-projects/07-humanoid/paper/codegen/tests/test\_xyz\_safety.py}.
Both files provide LLM credibility in generating tests and error
checks that the FDA and other regulators would expect in any human
factors review.

Subsection 8, Runtime Targets. Read every file in
\texttt{demo-projects/07-humanoid/paper/codegen/docker/} and the
\texttt{docker-compose.yml} at the codegen tree root. Cover MacOS,
Windows, and Linux runtime targets per the
\texttt{09-runtime-instructions/} directory.

Subsection 9, Execution. Read
\texttt{demo-projects/07-humanoid/paper/execution/README.md} and the
log files in
\texttt{demo-projects/07-humanoid/paper/execution/logs/}
(numbered 01 through 23). Walk through the schema ingest, smoke test,
cargo build, Rust runner, validation pass log progression so a
regulator can audit the run.

Required ASCII diagrams to include verbatim in this section:
\begin{itemize}
  \item \texttt{paper/codegen/diagrams/swarm\_dance.txt} as Figure 3.
  \item \texttt{paper/codegen/diagrams/ae\_response\_flow.txt} as
        Figure 5, placed at the end of subsection 9.
  \item \texttt{paper/execution/diagrams/01\_module\_dependency.txt}
        as Figure 6.
  \item \texttt{paper/execution/diagrams/03\_iteration\_sweep\_funnel.txt}
        as Figure 4 (subsection 6).
\end{itemize}

Required real-life feasibility insight to include in this section:
``Every Cartesian goal in \texttt{cartesian\_planner.py} is bounded by
the force budget caps and the inter robot minimum distance constants
named in the instructions README. These bounds are conservative
relative to what the H2 EDU spec sheet allows. The implication is
that a future regulator review can audit the safety envelope from
the constants alone, without re-running the simulation.''

Do not exceed 22 pages. Use single dashes only. Use \S\ for the
section sign. Avoid stranding 1 or 2 word lines on their own.]

The methods section will walk through the instruction layer, the 27
source files in \texttt{codegen/src/}, the schemas and configs, the
Cartesian coordinate command surface, the camaraderie behavior, the
32 iteration sweep, the test and error check files, the runtime
targets, and the execution log progression.

\sectiontable{Methods planning matrix}{%
Instructions & M1 & 1 & instructions tree & 10 subdir walkthrough & 1 to 2 pages & Names every subdir & Ready \\
Source files & M2 & 2 & codegen/src 27 files & Per file function & 4 to 5 pages & swarm\_coordinator at 191 & Ready \\
Schemas & M3 & 3 & schemas plus config & Schema by schema & 2 pages & Every JSON named & Ready \\
Cartesian & M4 & 4 & kinematics plus planner & x, y, z contract & 2 pages & Force caps named & Ready \\
Camaraderie & M5 & 5 & swarm\_coordinator.py & 191 line walkthrough & 2 pages & Figure 3 ASCII & Ready \\
Iteration & M6 & 6 & iterations/index.jsonl & 32 sweep summary & 1 to 2 pages & 32 across 4 sites & Ready \\
Tests & M7 & 7 & check\_errors plus tests & FDA credibility & 1 to 2 pages & Both files named & Ready \\
Runtime & M8 & 8 & docker plus compose & 3 OS targets & 1 to 2 pages & MacOS Win Linux & Ready \\
Execution & M9 & 9 & execution logs 01 to 23 & Run progression & 2 to 3 pages & Audit trail & Ready \\
}

\end{sectionaccent}


\section{Results}
\label{sec:results}
% ======================================================================
% sections/results.tex
%
% Results scaffold for the paper:
%   Triple Humanoid 24/7 Adverse Event Oncology Trial Response Team:
%   4-Site Rotation
%
% This file is a draft scaffold. A future Claude Code Opus 4.7 1M Max
% session reads the bracketed instruction blocks below, then reads the
% files those instructions name, then expands the scaffold into the
% finished results section. Target length is roughly 14 to 18 pages of
% the 70 plus page final paper.
% ======================================================================

\begin{sectionaccent}

[RESULTS INSTRUCTIONS FOR FUTURE CLAUDE CODE OPUS 4.7 1M MAX SESSION.
Write the results section as 7 subsections inside this single
\texttt{\textbackslash section\{Results\}} block. Every numeric claim must
come from a file named below. Do not introduce numbers that do not
appear in those files. Quote file paths in the prose so a reader can
verify each number.

Subsection 1, Per-Site LLM Decision Streams. Read the 4 site files in
\path{demo-projects/07-humanoid/paper/execution/data/site_runtime_runs/}:
\begin{itemize}
  \item \path{AT-01_llm_decisions.jsonl} (Atlanta, 62,330 bytes).
  \item \path{BO-01_llm_decisions.jsonl} (Boston, 62,330 bytes).
  \item \path{SD-01_llm_decisions.jsonl} (San Diego, 62,330 bytes).
  \item \path{SF-01_llm_decisions.jsonl} (San Francisco,
        62,330 bytes).
\end{itemize}
All 4 site files have the same byte length. Treat the byte length
match as a parallel-scope check: the 1 Hz broadcast cadence and the
shared schema produce decision logs of the same record count per
site. Read every line of \path{AT-01_llm_decisions.jsonl} (it is
the lead Atlanta file flagged as Hz data at scale for the project).
Then spot-check the other 3 site files for record count parity and
schema parity. If the records repeat across sites, call that out as a
duplication artifact in the runtime that production should fix
without retracting the parallel-scope claim. If the records differ
across sites despite identical byte length, call that out as
acceptable schema parity. Do this check explicitly and report the
finding in the prose.

Subsection 2, Iteration Sweep Aggregate. Read:
\begin{itemize}
  \item \path{demo-projects/07-humanoid/paper/codegen/data/iterations/index.jsonl}
        for the 32 iteration codegen index. State the 32 iteration
        count and the 4 site coverage.
  \item \path{demo-projects/07-humanoid/paper/execution/data/iterations_index_original.jsonl}
        and
        \path{demo-projects/07-humanoid/paper/execution/data/iterations_index_runtime.jsonl}
        for the original and runtime iteration indexes.
  \item \path{demo-projects/07-humanoid/paper/execution/data/iterations_aggregate.duckdb}
        as the DuckDB roll up. Discuss that DuckDB and JSONL coexist
        cleanly in the execution data layer. The DuckDB plus JSONL
        pairing demonstrates DuckDB proficiency for analytical queries
        and JSONL proficiency for append-only logs.
\end{itemize}

Subsection 3, Comparison Reports. Read:
\begin{itemize}
  \item \path{demo-projects/07-humanoid/paper/codegen/reports/comparison.json}.
  \item \path{demo-projects/07-humanoid/paper/codegen/reports/comparison_methodology.md}.
  \item \path{demo-projects/07-humanoid/paper/codegen/reports/report.md}.
  \item \path{demo-projects/07-humanoid/paper/codegen/reports/dashboard.html}.
        Treat the dashboard as the primary interpretability artifact
        for future paper viewers; the rendered HTML adds project
        interpretability without requiring the reader to inspect raw
        JSON. Embed a screenshot reference in the prose and link to
        the file path.
  \item \path{demo-projects/07-humanoid/paper/execution/reports/comparison_v0_4_0_copy.json}.
  \item \path{demo-projects/07-humanoid/paper/execution/reports/comparison_methodology_v0_4_0_copy.md}.
  \item \path{demo-projects/07-humanoid/paper/execution/reports/dashboard_v0_4_0_copy.html}.
  \item \path{demo-projects/07-humanoid/paper/execution/reports/execution_report_v0_5_0.md}.
  \item \path{demo-projects/07-humanoid/paper/execution/reports/report_v0_4_0_copy.md}.
\end{itemize}

Subsection 4, Comparison ASCII Visuals. Include verbatim inside
\texttt{asciibox} environments:
\begin{itemize}
  \item \path{demo-projects/07-humanoid/paper/execution/diagrams/04_comparison_radar.txt}
        as Figure 7.
  \item \path{demo-projects/07-humanoid/paper/execution/diagrams/01_module_dependency.txt}
        as Figure 8 if it was not already placed in the Methods
        section. If it was already placed there, only cross reference.
\end{itemize}

Subsection 5, Output Markdown Tables. Read:
\begin{itemize}
  \item \path{demo-projects/07-humanoid/paper/execution/outputs/comparison_summary_table.md}.
  \item \path{demo-projects/07-humanoid/paper/execution/outputs/iteration_sweep_table.md}.
  \item \path{demo-projects/07-humanoid/paper/execution/outputs/execution_tree.txt}.
  \item \path{demo-projects/07-humanoid/paper/execution/outputs/file_inventory.txt}.
\end{itemize}
Translate each markdown table into a \texttt{\textbackslash sectiontable}
LaTeX table using the 8 column shared macro. Do not change the data;
only re-render the format.

Subsection 6, Figures. Reference every figure in
\path{demo-projects/07-humanoid/paper/execution/figures/} (6 PNGs at
300 dpi: swarm architecture, response time histogram, camaraderie
heatmap, role rotation timeline, force budget, 4 site comparison).
Use \texttt{\textbackslash includegraphics} only if the figure file is
copied into the draft-paper folder as part of the future expansion.
Otherwise reference the file path in a caption note.

Subsection 7, Headline Metrics. Pull the headline numbers into one
final \texttt{\textbackslash sectiontable}:
\begin{itemize}
  \item Sites: 4 (SF-01, SD-01, BO-01, AT-01).
  \item Robots per site: 3 Unitree H2 EDU.
  \item Robots total: 12.
  \item LLM cadence: 1 Hz broadcast.
  \item Motion cadence: 10 Hz per robot.
  \item Monitoring window: 168 hours.
  \item AE volume: about 84 across the week.
  \item SAE volume: about 24 (CTCAE grade 3 or higher).
  \item Iterations: 32 deterministic sweeps.
  \item E-stop latency: 5 ms swarm wide.
  \item Per-arm cumulative force: 15 N during chest compressions,
        5 N otherwise.
  \item Inter-robot minimum distance: 1.2 m at rest, 0.4 m during
        shared task handoff.
  \item Cumulative cross-robot force: 22 N during 3-robot patient
        transfer.
  \item Camaraderie reduction factor: 3 in single robot error
        potential.
\end{itemize}
Every number in this final table must trace to a file in
\path{instructions/}, \texttt{codegen/}, or \texttt{execution/}.

Required real-life feasibility insight to include in this section:
``Identical byte length across the 4 site decision logs is not the
production end state. In a real 168 hour run the per-site record
counts will diverge because each site sees a different AE arrival
distribution. The current parity demonstrates that the runtime
stamps every site at 1 Hz on a shared schema; the next iteration is
to feed each site a site-specific arrival stream.''

Do not exceed 18 pages. Use single dashes only. Use \S\ for the
section sign. Avoid stranding 1 or 2 word lines on their own.]

The results section will present the per-site LLM decision streams,
the 32 iteration sweep aggregate, the comparison reports and ASCII
visuals, the markdown output tables converted to LaTeX tables, the
figures, and the final headline metrics table.

\sectiontable{Results planning matrix}{%
Decisions & R1 & 1 & runtime runs & 4 site walk & 2 pages & 62,330 bytes & Ready \\
Iterations & R2 & 2 & index, duckdb & 32 sweep & 2 pages & DuckDB JSONL & Ready \\
Reports & R3 & 3 & code, exec reports & Comparison & 3 pages & dashboard named & Ready \\
ASCII & R4 & 4 & exec diagrams & Figures 7, 8 & 1 to 2 pages & Verbatim text & Ready \\
Tables & R5 & 5 & exec outputs & Md to LaTeX & 2 to 3 pages & 8 col macro & Ready \\
Figures & R6 & 6 & exec figures & 6 PNG callouts & 1 to 2 pages & 300 dpi & Ready \\
Metrics & R7 & 7 & instructions metrics & Headline table & 1 page & Every number sourced & Ready \\
}

\end{sectionaccent}


\section{Discussion}
\label{sec:discussion}
% ======================================================================
% sections/discussion.tex
%
% Discussion scaffold for the paper:
%   Triple Humanoid 24/7 Adverse Event Oncology Trial Response Team:
%   4-Site Rotation
%
% This file is a draft scaffold. A future Claude Code Opus 4.7 1M Max
% session reads the bracketed instruction blocks below, then reads the
% files those instructions name, then expands the scaffold into the
% finished discussion. Target length is roughly 10 to 14 pages of the
% 70 plus page final paper.
% ======================================================================

\begin{sectionaccent}

[DISCUSSION INSTRUCTIONS FOR FUTURE CLAUDE CODE OPUS 4.7 1M MAX SESSION.
Write the discussion as 6 subsections inside this single
\texttt{\textbackslash section\{Discussion\}} block. Tie every claim back to a
results subsection.

Subsection 1, Thesis Reconciliation. Restate the thesis verbatim from
\texttt{demo-projects/07-humanoid/paper/instructions/README.md} and
walk through how each clause of the thesis is supported by:
\begin{itemize}
  \item ``On-premises repository based LLMs provide commands to
        humanoid robots'' is supported by
        \texttt{codegen/src/llm\_planner.py},
        \texttt{codegen/src/broadcaster.py}, and
        \texttt{codegen/src/comms\_intellectual.py}.
  \item ``based on real-time sensor data'' is supported by
        \texttt{codegen/src/sensor\_to\_xyz.py} and
        \texttt{codegen/src/peer\_state\_tracker.py}.
  \item ``controlled via x, y, z coordinates'' is supported by
        \texttt{codegen/src/cartesian\_planner.py},
        \texttt{codegen/src/kinematics.yaml}, and
        \texttt{codegen/src/xyz\_safety.py}.
  \item ``to administer synergistic treatment to patient adverse
        events'' is supported by
        \texttt{codegen/src/ctcae\_grader.py} and
        \texttt{codegen/src/physician\_escalation.py}.
  \item ``This workflow minimizes single robot error potential'' is
        supported by the 191 line
        \texttt{codegen/src/swarm\_coordinator.py} (the camaraderie
        factor of 3).
\end{itemize}

Subsection 2, Significance Of The On-Premises Repository LLM Pattern.
Discuss the quality and safety benefits of having one Claude Opus 4.7
1M instance per site rather than a remote cloud. Quality: every
inference traces to a known on-prem repository, which simplifies the
audit trail for the FDA RTCT pilot. Safety: a network partition
between the site and the public internet does not stop the 1 Hz
broadcast cadence. Speed: the 200 ms median latency named in the
instructions README is achievable on a local LAN. Read:
\begin{itemize}
  \item \texttt{demo-projects/07-humanoid/paper/instructions/README.md}
        Inventory at a Glance table.
  \item \texttt{demo-projects/07-humanoid/paper/codegen/src/llm\_planner.py}.
  \item \texttt{demo-projects/07-humanoid/paper/codegen/src/prompt\_frozen.md}.
\end{itemize}

Subsection 3, Significance Of The Cartesian Coordinate Surface. The
LLM does not author joint angles; it authors x, y, z goals and lets
the kinematics layer convert. This is the contract that lets a
regulator review the safety envelope from the constants alone. Read:
\begin{itemize}
  \item \texttt{demo-projects/07-humanoid/paper/codegen/src/kinematics.yaml}.
  \item \texttt{demo-projects/07-humanoid/paper/codegen/src/cartesian\_planner.py}.
  \item \texttt{demo-projects/07-humanoid/paper/codegen/src/xyz\_safety.py}.
  \item \texttt{demo-projects/07-humanoid/paper/codegen/tests/test\_xyz\_safety.py}.
\end{itemize}

Subsection 4, Significance Of The 32 Iteration Sweep And The
Competition Framing. Read:
\begin{itemize}
  \item \texttt{demo-projects/07-humanoid/paper/codegen/data/iterations/index.jsonl}
        (32 iterations across 4 sites and many columns).
  \item \texttt{demo-projects/07-humanoid/paper/codegen/src/iterate.py}
        and \texttt{demo-projects/07-humanoid/paper/codegen/src/compare\_agent.py}.
  \item Author's December 2025 paper
        \texttt{kawchak\_2025\_18029100}: ``Code Generation
        Competition: 16 Proprietary vs. Open-Source LLMs and Iterative
        Learning Based on FDA Adverse Event Reporting System.''
\end{itemize}
Tie the 32 iteration sweep back to the prior 16 LLM code generation
competition. The earlier competition established LLM effectiveness
against FDA adverse event data. The current 32 iteration sweep
extends that evidence to robot dispatch and Cartesian command code.

Subsection 5, FDA, CTCAE, And Other Regulators. Read:
\begin{itemize}
  \item \texttt{demo-projects/07-humanoid/paper/codegen/src/fda\_rtct\_submitter.py}
        (67 lines, HR 9505 SLA, ed25519 signing in production,
        placeholder in simulation).
  \item \texttt{demo-projects/07-humanoid/paper/codegen/src/ctcae\_grader.py}.
  \item \texttt{demo-projects/07-humanoid/paper/codegen/src/check\_errors.py}
        (LLM credibility in error checking code).
  \item \texttt{demo-projects/07-humanoid/paper/codegen/tests/test\_xyz\_safety.py}
        (LLM credibility in test code).
\end{itemize}
Write this subsection respectfully but advocate for accelerating
review timelines. The United States needs to remain Number 1 in the
world regarding patient safety, efficacy, and speed benefits in
oncological adverse event robot swarms for clinical trials. Frame the
review work as something the project has already checked off so the
agencies can focus on policy and not on rebuilding the audit trail.

Subsection 6, Prior Paper Context Integration. Read both Deep-Research
trees in
\texttt{demo-projects/07-humanoid/paper/inputs/Deep-Research-A/} and
\texttt{demo-projects/07-humanoid/paper/inputs/Deep-Research-B/}:
\begin{itemize}
  \item Deep-Research-A is the Medical Full Automation A source
        bundle: \texttt{README.md},
        \texttt{chunk1\_paper\_text.md},
        \texttt{chunk2\_tables\_references.md},
        \texttt{chunk3\_bibtex.md}.
  \item Deep-Research-B is the Medical Rapid Prototyping B source
        bundle: \texttt{chunk\_01\_executive\_summary\_and\_fully\_automated\_sponsor.md},
        \texttt{chunk\_02\_federated\_learning\_and\_patient\_prediction.md},
        \texttt{chunk\_03\_clinagent.md},
        \texttt{chunk\_04\_comparative\_overview\_and\_implications.md},
        \texttt{chunk\_05\_bibtex\_references.md}.
\end{itemize}
Compare and contrast the prior paper's Medical Full Automation A and
Medical Rapid Prototyping B framings with the current paper's 4 site
rotation framing. Position the current paper as the natural successor.

Required real-life feasibility insight to include in this section:
``The on-premises repository LLM pattern is feasible today on a single
high end server per site. The harder gap to close is not compute; it
is the IRB, sponsor, and site coordination required to enroll the
first real PAT-NET-001 cohort. The paper is a checklist of the
technical work that has been completed so that the human work of
governance can move forward without waiting on the technology.''

Do not exceed 14 pages. Use single dashes only. Use \S\ for the
section sign. Avoid stranding 1 or 2 word lines on their own.]

The discussion section will reconcile each clause of the thesis with
specific files, walk through the quality and safety significance of
the on-premises repository LLM pattern and the Cartesian coordinate
surface, position the 32 iteration sweep against the prior 16 LLM
competition, address FDA CTCAE and other regulators respectfully, and
integrate the prior paper context from Deep-Research-A and
Deep-Research-B.

\sectiontable{Discussion planning matrix}{%
Thesis & D1 & 1 & instructions README & Per clause source list & 2 pages & Every clause sourced & Ready \\
OnPrem & D2 & 2 & llm\_planner plus prompt & Quality, safety, speed & 2 pages & Audit trail framed & Ready \\
Cartesian & D3 & 3 & kinematics plus planner & Coordinate contract & 2 pages & Constants only review & Ready \\
Iteration & D4 & 4 & iterations index & 32 sweep significance & 2 pages & Cites kawchak\_2025 & Ready \\
Regulator & D5 & 5 & fda\_rtct plus ctcae & Respectful advocacy & 2 pages & HR 9505 named & Ready \\
PriorPaper & D6 & 6 & Deep-Research A and B & Successor framing & 2 to 3 pages & Both bundles named & Ready \\
}

\end{sectionaccent}


\section{Limitations and Future Work}
\label{sec:limitations}
% ======================================================================
% sections/limitations_future.tex
%
% Limitations and Future Work scaffold for the paper:
%   Triple Humanoid 24/7 Adverse Event Oncology Trial Response Team:
%   4-Site Rotation
%
% This file is a draft scaffold. A future Claude Code Opus 4.7 1M Max
% session reads the bracketed instruction blocks below, then reads the
% files those instructions name, then expands the scaffold into the
% finished section. Target length is roughly 6 to 8 pages of the 70
% plus page final paper.
% ======================================================================

\begin{sectionaccent}

[LIMITATIONS AND FUTURE WORK INSTRUCTIONS FOR FUTURE CLAUDE CODE OPUS
4.7 1M MAX SESSION. Write 5 subsections inside this single
\texttt{\textbackslash section} block. Every limitation must trace to a
file so the reader can verify the bound and so the future work
roadmap stays anchored to the real artifact.

Subsection 1, Simulation Versus Real Robot Limitations. Read:
\begin{itemize}
  \item \path{demo-projects/07-humanoid/paper/codegen/src/fda_rtct_submitter.py}
        (production uses ed25519 signing, simulation uses a placeholder
        signature).
  \item \path{demo-projects/07-humanoid/paper/codegen/src/robot_loop.cpp}
        (the in-robot 10 Hz loop runs in a simulator harness, not on
        the H2 EDU body).
  \item \path{demo-projects/07-humanoid/paper/codegen/src/sensor_to_xyz.py}
        (sensor frames are synthetic in the current run).
\end{itemize}
State that no real Unitree H2 EDU body executed any of the broadcast
commands in this paper. State that the FDA RTCT pilot intake API is
called only in simulation. State that the patient identifiers are
synthetic and have the form \texttt{PAT-NET-001-PNNN}.

Subsection 2, Per-Site Data Identity Limitation. Read the 4 site
files in
\path{demo-projects/07-humanoid/paper/execution/data/site_runtime_runs/}.
All 4 files share the same byte length (62,330 bytes each). State
that this parity is acceptable as a schema check and is not yet a
production behavior. State that the next iteration should feed each
site a site-specific AE arrival stream so the record content
diverges by site.

Subsection 3, Excluded Inputs. Read the Notes section of
\path{demo-projects/07-humanoid/paper/instructions/README.md}. The
original Prompt 07 extra hours dataset (hour 56 through hour 83 under
\path{physical-ai-oncology-trials/new-trial/national-24-7-trial/extra-hours/})
is excluded from the current input scope. Only hour 00 through hour
55 are read. State this explicitly and identify the missing hours as
a future inclusion roadmap item.

Subsection 4, Missing Repository Update Instruction Directory. Read
the Repository Structure block of
\path{demo-projects/07-humanoid/paper/instructions/README.md}. The
\path{10-repository-update-instructions/} subdirectory is named in
the README but is not yet present on disk. State this explicitly as a
documented limitation and as the next commit for the instruction
tree.

Subsection 5, Future Work Roadmap. Read:
\begin{itemize}
  \item \path{demo-projects/07-humanoid/paper/inputs/README.md}
        for the prior paper's future work framing.
  \item \path{demo-projects/07-humanoid/paper/inputs/Deep-Research-A/chunk2_tables_references.md}
        for the Medical Full Automation A roadmap items.
  \item \path{demo-projects/07-humanoid/paper/inputs/Deep-Research-B/chunk_04_comparative_overview_and_implications.md}
        for the Medical Rapid Prototyping B roadmap items.
  \item Author's December 2025 paper
        \path{kawchak_2025_18029100} for the prior limitations
        that this paper carries forward.
\end{itemize}
Roadmap items to enumerate explicitly:
\begin{itemize}
  \item Move from simulator harness to a real Unitree H2 EDU body
        with the kinematics in
        \path{codegen/src/kinematics.yaml} as the live target.
  \item Replace the placeholder signature in
        \path{codegen/src/fda_rtct_submitter.py} with the real
        ed25519 signing chain.
  \item Enroll a real PAT-NET-001 cohort under an IRB protocol and a
        signed sponsor agreement.
  \item Add the
        \path{instructions/10-repository-update-instructions/}
        directory.
  \item Extend the AE arrival stream so the 4 site decision logs
        diverge in content, not only in byte length.
  \item Include the hour 56 through hour 83 extra hours dataset from
        \texttt{physical-ai-oncology-trials} once the supporting
        analysis is ready.
  \item Upgrade the per-site compute path to NVIDIA Jetson AGX Thor
        2070 TOPS where the workload demands it.
\end{itemize}

Required real-life feasibility insight to include in this section:
``The largest gap between this paper and a real Phase 2 trial response
team is not technology. It is the IRB, sponsor, and FDA RTCT pilot
governance. The technology layer is feasible today. The governance
layer is the work that the paper invites the agencies, sponsors, and
sites to take on.''

Do not exceed 8 pages. Use single dashes only. Use \S\ for the
section sign. Avoid stranding 1 or 2 word lines on their own.]

The limitations and future work section will name the simulator
versus real robot gap, the per-site data identity limitation, the
excluded extra hours dataset, the missing repository update
instruction directory, and the future work roadmap including the IRB
and FDA RTCT pilot enrollment path.

\sectiontable{Limitations and Future Work planning matrix}{%
SimGap & L1 & 1 & robot, submit. & Sim vs real & 1 to 2 pages & No real H2 run & Ready \\
Parity & L2 & 2 & runtime runs & Identity check & 1 page & 62,330 bytes & Ready \\
Excluded & L3 & 3 & instr Notes & Hour 56 to 83 & 1 page & By README & Ready \\
Missing & L4 & 4 & instr structure & 10-repo gap & 1 page & Not present & Ready \\
Future & L5 & 5 & inputs, DR-A, DR-B & Roadmap list & 2 to 3 pages & 7 items & Ready \\
}

\end{sectionaccent}


\section{Conclusions}
\label{sec:conclusions}
% ======================================================================
% sections/conclusions.tex
%
% Conclusions scaffold for the paper:
%   Triple Humanoid 24/7 Adverse Event Oncology Trial Response Team:
%   4-Site Rotation
%
% This file is a draft scaffold. A future Claude Code Opus 4.7 1M Max
% session reads the bracketed instruction blocks below, then reads the
% files those instructions name, then expands the scaffold into the
% finished conclusions section. Target length is roughly 3 to 5 pages
% of the 70 plus page final paper.
% ======================================================================

\begin{sectionaccent}

[CONCLUSIONS INSTRUCTIONS FOR FUTURE CLAUDE CODE OPUS 4.7 1M MAX
SESSION. Write 4 short subsections inside this single
\texttt{\textbackslash section} block. Conclusions does not introduce new
files. Every file named here must also be cited earlier in the paper.

Subsection 1, Thesis Restatement. Restate the thesis verbatim. Do not
paraphrase. The thesis appears in
\texttt{demo-projects/07-humanoid/paper/instructions/README.md} under
the heading ``Thesis''.

Subsection 2, Evidence Roll Up. List the 10 evidence anchors that
support the thesis. Each anchor is one bullet, one sentence:
\begin{itemize}
  \item \texttt{codegen/src/swarm\_coordinator.py} (191 lines)
        implements the camaraderie logic that animates the 3 robot
        Unitree H2 EDU camarade swarm at one site. The compact
        line count is itself evidence of clean coordination logic.
  \item \texttt{codegen/src/fda\_rtct\_submitter.py} (67 lines)
        submits CTCAE grade 2 or higher records to the FDA RTCT
        pilot intake API within 1 hour per the HR 9505 SLA, using
        ed25519 signing in production.
  \item The 27 files in \texttt{codegen/src/} span Python, C++,
        Rust, YAML, and Markdown and together implement the full
        sensor to LLM to robot pipeline.
  \item \texttt{codegen/data/iterations/index.jsonl} carries 32
        deterministic iteration sweeps across the 4 trial sites and
        many columns.
  \item \texttt{codegen/reports/dashboard.html} adds project
        interpretability for future paper readers.
  \item \texttt{codegen/src/check\_errors.py} provides LLM
        credibility in generating error checking code.
  \item \texttt{codegen/tests/test\_xyz\_safety.py} provides LLM
        credibility in generating test code.
  \item \texttt{codegen/src/h2\_dispatcher.py} provides per-site H2
        dispatch proficiency.
  \item \texttt{execution/data/site\_runtime\_runs/AT-01\_llm\_decisions.jsonl}
        plus the matching SF-01, SD-01, and BO-01 files show 1 Hz
        decision logs at scale for all 4 sites.
  \item \texttt{execution/data/iterations\_aggregate.duckdb} and the
        JSONL files in \texttt{execution/data/} together
        demonstrate DuckDB and JSONL proficiency in the same data
        layer.
\end{itemize}

Subsection 3, Industry Checklist Framing. Present the entire body of
work as a sequence of checklist items that the industry no longer
has to redo. Pull the checklist from the Inventory at a Glance table
in
\texttt{demo-projects/07-humanoid/paper/instructions/README.md}. Mark
each item as done so the next team can build on top of it. The point
of the framing is to save the industry the cost of repeating the
groundwork.

Subsection 4, Forward Statement. Close with the United States
patient safety, efficacy, and speed framing. The United States needs
to remain Number 1 in the world regarding patient safety, efficacy,
and speed benefits in oncological adverse event robot swarms for
clinical trials. Frame the closing sentence as an invitation to the
FDA, sponsors, sites, and IRBs to advance the new field of Physical
AI oncology trial surgical procedures together. Cite the author's
December 2025 paper
\texttt{kawchak\_2025\_18029100} one final time as the prior
evidence that LLMs are effective on FDA adverse event data.

Required real-life feasibility insight to include in this section:
``The pipeline in this paper is small enough to audit, large enough
to deploy at 4 sites, and tight enough to honor the FDA RTCT 1 hour
SLA. The work is ready for the human governance step.''

Do not exceed 5 pages. Use single dashes only. Use \S\ for the
section sign. Avoid stranding 1 or 2 word lines on their own.]

The conclusions section will restate the thesis verbatim, roll up the
10 evidence anchors, frame the body of work as an industry checklist
of completed steps, and close with the United States patient safety,
efficacy, and speed forward statement plus the final citation to
\texttt{kawchak\_2025\_18029100}.

\sectiontable{Conclusions planning matrix}{%
Thesis & C1 & 1 & instructions README & Verbatim restatement & 1 paragraph & No paraphrase & Ready \\
Evidence & C2 & 2 & 10 anchors & Bullet roll up & 1 page & 10 bullets & Ready \\
Checklist & C3 & 3 & Inventory at a Glance & Done checklist framing & 1 to 2 pages & Industry handoff & Ready \\
Forward & C4 & 4 & kawchak\_2025\_18029100 & US leadership close & 1 paragraph & Final citation & Ready \\
}

\end{sectionaccent}


\phantomsection
\addcontentsline{toc}{section}{References}
% \nocite{*} forces every entry in references.bib to appear in the
% printed bibliography even when the scaffold prose does not cite it.
% The future Claude Code Opus 4.7 1M Max session that expands the
% scaffold can replace this with explicit \cite commands once the
% paper body cites each reference at least once. Until then, this
% line guarantees every doi and url in references.bib shows up.
\nocite{*}
\bibliographystyle{ieeetr}
\bibliography{references}

% ======================================================================
% sections/back_matter.tex
%
% Back matter for the paper:
%   Triple Humanoid 24/7 Adverse Event Oncology Trial Response Team:
%   4-Site Rotation
%
% This file holds the Acknowledgments, Ethical Disclosures, Rights and
% Permissions, Cite This Article, and Data Availability sections, in
% that order. The text below is final-form text. A future Claude Code
% Opus 4.7 1M Max session should keep this content and only update the
% DOI placeholder when the Zenodo record is minted.
% ======================================================================

\phantomsection
\addcontentsline{toc}{section}{Acknowledgments}
\section*{Acknowledgments}

The author would like to acknowledge Anthropic for providing access to
Claude Code for instruction, code, execution, along with the paper
generation process; OpenAI for providing access to ChatGPT for deep
research summaries, and Google Gemini AI Overview for additional
search assistance.

\phantomsection
\addcontentsline{toc}{section}{Ethical Disclosures}
\section*{Ethical Disclosures}

The author of the article declares no competing interests.

\phantomsection
\addcontentsline{toc}{section}{Rights and Permissions}
\section*{Rights and Permissions}

This article is distributed under the terms of the Creative Commons
Attribution 4.0 International License (CC BY 4.0), which permits
unrestricted use, distribution, and reproduction in any medium,
provided the original author(s) and source are properly credited, a
link to the Creative Commons license is provided, and any
modifications made are indicated. To view a copy of this license,
visit \href{https://creativecommons.org/licenses/by/4.0/}{https://creativecommons.org/licenses/by/4.0/}.

\phantomsection
\addcontentsline{toc}{section}{Cite This Article}
\section*{Cite This Article}

Kawchak K. Triple Humanoid 24/7 Adverse Event Oncology Trial Response
Team: 4-Site Rotation. Zenodo. 2026; \href{https://doi.org/10.5281/zenodo.xxxxxxxx}{10.5281/zenodo.xxxxxxxx}.

\phantomsection
\addcontentsline{toc}{section}{Data Availability}
\section*{Data Availability}

[DATA AVAILABILITY INSTRUCTIONS FOR FUTURE CLAUDE CODE OPUS 4.7 1M
MAX SESSION. Write 1 paragraph that points the reader at the exact
repository directories where the data, code, schemas, configs, and
reports live. Use the file paths below. Do not omit any of them.

\begin{itemize}
  \item Source code: \path{demo-projects/07-humanoid/paper/codegen/src/}
        (27 files spanning Python, C++, Rust, YAML, and Markdown).
  \item Tests: \path{demo-projects/07-humanoid/paper/codegen/tests/}
        (notably \path{test_xyz_safety.py}).
  \item Schemas: \path{demo-projects/07-humanoid/paper/codegen/schemas/}.
  \item Configs: \path{demo-projects/07-humanoid/paper/codegen/config/}.
  \item Data: \path{demo-projects/07-humanoid/paper/codegen/data/}
        and \path{demo-projects/07-humanoid/paper/execution/data/}.
        The iteration index lives at
        \path{codegen/data/iterations/index.jsonl}. The DuckDB
        aggregate lives at
        \path{execution/data/iterations_aggregate.duckdb}. The
        per-site decision streams live in
        \path{execution/data/site_runtime_runs/} as
        \path{AT-01_llm_decisions.jsonl},
        \path{BO-01_llm_decisions.jsonl},
        \path{SD-01_llm_decisions.jsonl}, and
        \path{SF-01_llm_decisions.jsonl}.
  \item Reports: \path{demo-projects/07-humanoid/paper/codegen/reports/}
        and \path{demo-projects/07-humanoid/paper/execution/reports/}.
        The interpretability dashboard lives at
        \path{codegen/reports/dashboard.html}.
  \item ASCII diagrams:
        \path{demo-projects/07-humanoid/paper/codegen/diagrams/}
        and
        \path{demo-projects/07-humanoid/paper/execution/diagrams/}.
  \item Figures: \path{demo-projects/07-humanoid/paper/execution/figures/}
        (6 PNGs at 300 dpi).
  \item Instructions tree:
        \path{demo-projects/07-humanoid/paper/instructions/}
        with subdirectories \path{00-project-overview/} through
        \path{09-runtime-instructions/}.
  \item Prior paper inputs:
        \path{demo-projects/07-humanoid/paper/inputs/} including
        \path{Deep-Research-A/} and \path{Deep-Research-B/}.
  \item Companion repository:
        \path{https://github.com/kevinkawchak/physical-ai-oncology-trials},
        with hours 00 through 55 of
        \path{new-trial/national-24-7-trial/} included as input
        and hours 56 through 83 excluded per the instructions
        README.
\end{itemize}

Close the paragraph by linking to the project DOI placeholder
\path{https://doi.org/10.5281/zenodo.xxxxxxxx} and to the
instructions DOI \path{https://doi.org/10.5281/zenodo.18445179}.]

The data availability statement will name every codegen, execution,
instruction, and input directory so the reader can navigate to the
exact source of every claim in the paper, and will link the project
DOI placeholder and the instructions DOI in a single closing
sentence.

\sectiontable{Back matter disclosure matrix}{%
Acks & B1 & 1 & Anthr, OAI, Goog & Record & Ack block & AI disclosure & Ready \\
Ethics & B2 & 2 & Author & No competing interest & Ethical block & Single sentence & Ready \\
Rights & B3 & 3 & CC BY 4.0 & License terms & Permissions block & License link & Ready \\
Cite & B4 & 4 & DOI placeholder & Suggested citation & Cite block & DOI link & Ready \\
Data & B5 & 5 & Whole paper tree & Per directory pointer & Data block & 9 dirs named & Ready \\
}

\end{document}
